
\section*{Introduction}
\addcontentsline{toc}{section}{Introduction}

The healthcare sector generates vast quantities of heterogeneous data, from high-resolution medical images and longitudinal electronic health records to genomic sequences and clinical narratives. While deep learning has achieved remarkable success across these domains, the conventional paradigm of training monolithic, dense models faces fundamental scalability constraints: increasing model capacity uniformly activates all parameters for every input, resulting in prohibitive computational costs that limit deployment in resource-constrained clinical environments \cite{kaplan2020scaling}. The Mixture of Experts (MoE) architecture offers an elegant solution to this dilemma by enabling \textit{conditional computation}, where only a subset of model parameters is activated for each input, thereby decoupling total model capacity from inference cost \cite{jacobs1991adaptive, shazeer2017outrageously}.

The core insight of MoE aligns naturally with the structure of medical data and clinical reasoning. Diseases manifest differently across patients, imaging modalities capture complementary anatomical information, and clinical decision-making inherently involves consulting domain-specific knowledge. An MoE framework mirrors this structure by routing each input to specialized expert sub-networks, each potentially capturing distinct patterns corresponding to disease subtypes, tissue morphologies, imaging artifacts, or patient demographics \cite{masoudnia2014mixture}. This conditional activation paradigm has proven especially valuable in healthcare, where data distributions are long-tailed, class imbalances are severe, and the cost of errors can be life-threatening.

Recent years have witnessed an explosive growth in MoE-based approaches for healthcare applications. In medical imaging, MoE architectures have been deployed to handle variable image quality, multi-site heterogeneity, and multi-task segmentation \cite{rahman2025personalized_moe, xie2024moe_pathology}. In clinical natural language processing, MoE layers have been integrated into large language models (LLMs) to specialize different experts for distinct medical sub-domains, improving diagnostic reasoning while reducing hallucination \cite{bechar2025biocancer_moe, chen2026cor_mole}. In drug discovery, MoE frameworks have enabled joint modeling of diverse molecular properties and multi-scale biological representations \cite{acharya2025mobius}. In electronic health record (EHR) analysis, MoE-based multi-task learning has simultaneously predicted multiple clinical outcomes by routing patient records through task-specific experts \cite{rajkomar2018ehr_deep}.

This Review provides a comprehensive examination of MoE architectures applied to healthcare. We begin by establishing the mathematical and architectural foundations of MoE, including gating mechanisms, expert specialization strategies, and load balancing. We then systematically survey applications across five major healthcare domains: medical imaging and computational pathology, clinical NLP and medical LLMs, drug discovery and genomics, EHR analysis, and multi-modal clinical decision support. We critically evaluate the unique challenges that arise when deploying MoE in clinical settings, including interpretability requirements, regulatory considerations, expert collapse, and data privacy constraints. We also outline future directions where MoE architectures may fundamentally reshape clinical AI, including personalized routing, federated expert training, and integration into healthcare foundation models.
